\documentclass[12pt]{article}
\usepackage[T1]{fontenc}
\usepackage[utf8]{inputenc}
\usepackage{lmodern}
\usepackage{textcomp}
\usepackage{enumitem}
\usepackage{geometry}
\geometry{margin=1in}
\begin{document}
\begin{center}
Answer Key - Philosophy - Undergraduate Level Assessment 20 \
\small (Answers are ordered exactly as in the question paper.)
\end{center}

\section*{Multiple Choice}
\begin{enumerate}[label=\arabic*.]
\item \textbf{Answer:} A
\\ \textit{Options:} A) a pupa.; B) an embryo.; C) a gamete.; D) a larva.; E) a neonate.
\\ \textit{Note:} The correct answer is a pupa because it refers to the life stage of certain insects during which they undergo transformation into their adult form, exhibiting a distinct and characteristic structure of their species.
\item \textbf{Answer:} D
\\ \textit{Options:} A) imaginary; B) interdependent; C) both nonexistent; D) one and the same thing; E) two distinct things
\\ \textit{Note:} Berkeley's philosophy of idealism posits that objects only exist as they are perceived, leading him to conclude that to exist and to be perceived are inseparably linked, making them "one and the same thing."
\item \textbf{Answer:} A
\\ \textit{Options:} A) The Battle of Badr; B) The Battle of Yamama; C) The Battle of Hunayn; D) The Battle of Mut'ah; E) The Battle of Tabuk
\\ \textit{Note:} The Battle of Badr, fought in 624 CE, was a pivotal victory for the early Muslim community against the Quraysh of Mecca, which not only solidified Muhammad's leadership and the legitimacy of Islam but also set the stage for subsequent events that ultimately led to the peaceful conquest and conversion of Mecca to Islam in 630 CE.
\item \textbf{Answer:} A
\\ \textit{Options:} A) Kantian distributionism.; B) libertarian cosmopolitanism.; C) Rawlsian justice.; D) Communitarianism.; E) Anarchist cosmopolitanism.
\\ \textit{Note:} Carens advocates for a Kantian distributionism because he emphasizes the moral imperative of distributing resources and opportunities based on principles of justice and equality, reflecting Kant's emphasis on human dignity and ethical treatment.
\item \textbf{Answer:} A
\\ \textit{Options:} A) love.; B) truth.; C) justice.; D) faith.; E) greed.
\\ \textit{Note:} Augustine argues that evil is a privation of good, and since love is fundamentally good, its absence creates a void where evil can manifest, indicating that love is essential for the existence of moral value and, consequently, for the existence of evil.
\end{enumerate}

\section*{True / False}
\begin{enumerate}[label=\arabic*.]
\setcounter{enumi}{5}
\item \textbf{Answer:} True
\\ \textit{Options:} A) True; B) False
\\ \textit{Note:} Ghosa, Apala, and Lopamurda are acknowledged in the Rigveda as notable female poets, highlighting the presence and significance of women's voices in the early Vedic literature.
\item \textbf{Answer:} False
\\ \textit{Options:} A) True; B) False
\\ \textit{Note:} Murtis are not considered 'spirits' but rather are physical representations or embodiments of deities in Hinduism, serving as focal points for worship and devotion rather than being viewed as spiritual entities themselves.
\item \textbf{Answer:} False
\\ \textit{Options:} A) True; B) False
\\ \textit{Note:} The ambrosial hours for Sikhs, known as "Amrit Vela," are actually observed from approximately 3 a.m. to 6 a.m., not from 6 to 9 a.m.
\item \textbf{Answer:} True
\\ \textit{Options:} A) True; B) False
\\ \textit{Note:} The Old Babylonian version of Gilgamesh underwent modifications into the standard version during the late second millennium BCE as part of a broader effort to unify and standardize literary and cultural texts in ancient Mesopotamia.
\item \textbf{Answer:} True
\\ \textit{Options:} A) True; B) False
\\ \textit{Note:} Baxter asserts that introducing new species can enhance biodiversity and improve ecosystem resilience, thereby addressing environmental challenges and fostering a more balanced ecological system.
\end{enumerate}

\section*{Long Form Response}
\begin{enumerate}[label=\arabic*.]
\setcounter{enumi}{10}
\item \textbf{Answer:} Michael Walzer identifies four main justifications for terrorism: the struggle against oppression, the defense of a community, the pursuit of political goals, and the response to injustices. While these justifications may suggest a context in which terrorism could be seen as a legitimate form of resistance, they do not support the notion of terrorism as a form of freedom of speech. Walzer argues that freedom of speech entails dialogue and reasoned debate, whereas terrorism employs violence and intimidation, undermining the possibility of constructive discourse. Thus, framing terrorism as an expression of free speech fails to align with Walzer's ethical considerations, which prioritize the moral implications of actions over mere expression. In this way, terrorism cannot be reconciled with the principles of freedom of speech that encourage mutual respect and understanding among individuals.
\\ \textit{Note:} correct because it emphasizes that while Walzer's justifications for terrorism relate to resistance and community defense, they fundamentally contradict the principles of freedom of speech, which relies on non-violent communication and rational discourse rather than violent acts.
\item \textbf{Answer:} Socrates objects to Euthyphro's definition of the holy as merely the act of prosecuting wrongdoers by arguing that this definition is overly narrow and fails to capture the true essence of holiness. He points out that if prosecuting wrongdoers is the sole criterion for holiness, then many other actions, beliefs, or qualities could also be considered holy, which dilutes the specificity of the term. For instance, one could argue that acts of charity or piety might also be holy, suggesting that holiness encompasses a broader range of virtuous actions. This critique leads to the implication that holiness cannot simply be defined by a single action or characteristic; instead, it requires a more comprehensive understanding that accounts for various forms of moral and ethical behavior. Thus, Socrates highlights the complexity of holiness, urging a deeper examination of what it truly means to be holy beyond mere legalistic definitions.
\\ \textit{Note:} correct because Socrates critiques Euthyphro's definition as insufficiently comprehensive, arguing that it reduces holiness to a singular action-prosecuting wrongdoers-while neglecting other essential qualities or actions that may also embody holiness, thereby challenging the completeness of the definition.
\end{enumerate}

\vfill
\noindent\textit{End of Answer Key}
\end{document}